\chapter{2007年陕西卷（文）}

\section{单选题}

\begin{problem}
已知全集 \(U = \{1, 2, 3, 4, 5, 6\}\)，集合 \(A = \{2, 3, 6\}\)，则集合 \(\complement_U A\) 等于（\quad）

\choices
    {\(\{1, 4\}\)}
    {\(\{4, 5\}\)}
    {\(\{1, 4, 5\}\)}
    {\(\{2, 3, 6\}\)}
\end{problem}

\begin{answer}
C
\end{answer}

\begin{solution}
补集由全集中不属于集合 \(A\) 的元素组成，因此 \(\complement_U A=\{1,4,5\}\)，故选 C.
\end{solution}

\begin{problem}
函数 \(f(x) = \lg \sqrt{1 - x^2}\) 的定义域为（\quad）

\choices
    {\([0, 1]\)}
    {\((-1, 1)\)}
    {\([-1, 1]\)}
    {\((-\infty, -1) \cup (1, +\infty)\)}
\end{problem}

\begin{answer}
B
\end{answer}

\begin{solution}
因为对数的真数必须为正，所以 \(\sqrt{1-x^2}>0\)，即 \(1-x^2>0\)，解得 \(-1<x<1\). 因此定义域为 \((-1,1)\)，故选 B.
\end{solution}

\begin{problem}
抛物线 \(x^2 = y\) 的准线方程是（\quad）

\choices
    {\(4x + 1 = 0\)}
    {\(4y + 1 = 0\)}
    {\(2x + 1 = 0\)}
    {\(2y + 1 = 0\)}
\end{problem}

\begin{answer}
B
\end{answer}

\begin{solution}
将抛物线方程写成标准形式 \(x^2=4py\)，比较系数得 \(4p=1\)，所以 \(p=\frac14\). 对于抛物线 \(x^2=4py\)，准线方程为 \(y=-p\)，即 \(4y+1=0\)，故选 B.
\end{solution}

\begin{problem}
已知 \(\sin \alpha = \frac{\sqrt{5}}{5}\)，则 \(\sin^4 \alpha - \cos^4 \alpha\) 的值为（\quad）

\choices
    {\(-\frac{3}{5}\)}
    {\(-\frac{1}{5}\)}
    {\(\frac{1}{5}\)}
    {\(\frac{3}{5}\)}
\end{problem}

\begin{answer}
A
\end{answer}

\begin{solution}
由 \(\sin\alpha=\frac{\sqrt5}{5}\) 得 \(\sin^2\alpha=\frac15\)，从而 \(\cos^2\alpha=1-\sin^2\alpha=\frac45\). 因此 \(\sin^4\alpha-\cos^4\alpha=(\sin^2\alpha-\cos^2\alpha)(\sin^2\alpha+\cos^2\alpha)=\frac15-\frac45=-\frac35\)，故选 A.
\end{solution}

\begin{problem}
等差数列 \(\{a_n\}\) 的前 \(n\) 项和为 \(S_n\)，若 \(S_2 = 2\)，\(S_4 = 10\)，则 \(S_6\) 等于（\quad）

\choices
    {\(12\)}
    {\(18\)}
    {\(24\)}
    {\(42\)}
\end{problem}

\begin{answer}
C
\end{answer}

\begin{solution}
设等差数列的首项为 \(a_1\)，公差为 \(d\). 由前 \(n\) 项和公式，
\(S_2=2a_1+d=2\)，\(S_4=4a_1+6d=10\). 将第一式乘以 \(2\) 得 \(4a_1+2d=4\)，用第二式减去该式得 \(4d=6\)，所以 \(d=\frac32\)，再由 \(2a_1+d=2\) 得 \(a_1=\frac14\). 因而
\(S_6=\frac{6}{2}\left(2a_1+5d\right)=3\left(\frac12+\frac{15}{2}\right)=24\)，故选 C.
\end{solution}

\begin{problem}
某商场有四类食品，其中粮食类，植物油类，动物性食品类以及果蔬类分别有 \(40\) 种，\(10\) 种，\(30\) 种，\(20\) 种，现从中抽取一个容量为 \(20\) 的样本进行食品安全检测. 若采用分层抽样的方法抽取样本，则抽取的植物油类与果蔬类食品种数之和是（\quad）

\choices
    {\(4\)}
    {\(5\)}
    {\(6\)}
    {\(7\)}
\end{problem}

\begin{answer}
C
\end{answer}

\begin{solution}
四类食品总数为 \(40+10+30+20=100\)，抽取容量为 \(20\)，所以分层抽样的抽样比为 \(\frac{20}{100}=\frac15\). 植物油类抽取 \(10\times\frac15=2\) 种，果蔬类抽取 \(20\times\frac15=4\) 种，二者之和为 \(2+4=6\)，故选 C.
\end{solution}

\begin{problem}
\(\mathrm{Rt}\triangle ABC\) 的三个顶点在半径为 \(13\) 的球面上，直角边的长分别为 \(6\) 和 \(8\)，则球心到平面 \(ABC\) 的距离是（\quad）

\choices
    {\(5\)}
    {\(6\)}
    {\(10\)}
    {\(12\)}
\end{problem}

\begin{answer}
D
\end{answer}

\begin{solution}
直角三角形的斜边长为 \(\sqrt{6^2+8^2}=10\)，所以其外接圆半径为 \(\frac{10}{2}=5\). 该三角形的外接圆是球面与平面 \(ABC\) 的截圆，设球心到平面 \(ABC\) 的距离为 \(h\)，则
\(h^2+5^2=13^2\)，从而 \(h=\sqrt{169-25}=12\)，故选 D.
\end{solution}

\begin{problem}
设函数 \(f(x) = 2^x + 1\)（\(x \in \R\)）的反函数为 \(f^{-1}(x)\)，则函数 \(y = f^{-1}(x)\) 的图象是（\quad）

\choices
    {\choicebitmap[width=0.15\paperwidth]{img/2007/shaanxi_liberal/q8_optA_clean.png}}
    {\choicebitmap[width=0.15\paperwidth]{img/2007/shaanxi_liberal/q8_optB_clean.png}}
    {\choicebitmap[width=0.15\paperwidth]{img/2007/shaanxi_liberal/q8_optC_clean.png}}
    {\choicebitmap[width=0.15\paperwidth]{img/2007/shaanxi_liberal/q8_optD_clean.png}}
\end{problem}

\begin{answer}
A
\end{answer}

\begin{solution}
由 \(y=2^x+1\) 得 \(x=\log_2(y-1)\)，所以反函数为 \(y=\log_2(x-1)\)，其定义域为 \(x>1\). 该图象的渐近线为 \(x=1\)，并且经过点 \((2,0)\)，与选项 A 的图象相符，故选 A.
\end{solution}

\begin{problem}
已知双曲线 \(C: \frac{x^2}{a^2} - \frac{y^2}{b^2} = 1\)（\(a > 0\)，\(b > 0\)），以 \(C\) 的右焦点为圆心且与 \(C\) 的渐近线相切的圆的半径是（\quad）

\choices
    {\(a\)}
    {\(b\)}
    {\(\sqrt{ab}\)}
    {\(\sqrt{a^2 + b^2}\)}
\end{problem}

\begin{answer}
B
\end{answer}

\begin{solution}
双曲线的右焦点为 \((c,0)\)，其中 \(c=\sqrt{a^2+b^2}\)，其渐近线可写为 \(bx-ay=0\) 或 \(bx+ay=0\). 点 \((c,0)\) 到任一条渐近线的距离为
\(\frac{bc}{\sqrt{a^2+b^2}}=\frac{bc}{c}=b\). 因而所求圆的半径为 \(b\)，故选 B.
\end{solution}

\begin{problem}
已知 \(P\) 为平面 \(\alpha\) 外一点，直线 \(l \subset \alpha\)，点 \(Q \in l\)，记点 \(P\) 到平面 \(\alpha\) 的距离为 \(a\)，点 \(P\) 到直线 \(l\) 的距离为 \(b\)，点 \(P\)，\(Q\) 之间的距离为 \(c\)，则（\quad）

\choices
    {\(a \leqslant b \leqslant c\)}
    {\(c \leqslant a \leqslant b\)}
    {\(a \leqslant c \leqslant b\)}
    {\(b \leqslant c \leqslant a\)}
\end{problem}

\begin{answer}
A
\end{answer}

\begin{solution}
设 \(H\) 为点 \(P\) 在平面 \(\alpha\) 上的垂足，则 \(PH=a\). 因为直线 \(l\subset\alpha\)，对任意点 \(X\in l\)，均有 \(PX\geqslant PH=a\)，所以点 \(P\) 到直线 \(l\) 的距离满足 \(b\geqslant a\). 又 \(Q\in l\)，故 \(c=PQ\geqslant b\). 因此 \(a\leqslant b\leqslant c\)，故选 A.
\end{solution}

\begin{problem}
给出如下三个命题：

\begin{circlelist}
\item 设 \(a\)，\(b \in \R\)，且 \(ab \neq 0\)，若 \(\frac{b}{a} > 1\)，则 \(\frac{a}{b} < 1\)；
\item 四个非零实数 \(a\)，\(b\)，\(c\)，\(d\) 依次成等比数列的充要条件是 \(ad = bc\)；
\item 若 \(f(x) = \log_2 x\)，则 \(f(|x|)\) 是偶函数.
\end{circlelist}

其中正确命题的序号是（\quad）

\choices
    {\(\circled{1}\circled{2}\)}
    {\(\circled{2}\circled{3}\)}
    {\(\circled{1}\circled{3}\)}
    {\(\circled{1}\circled{2}\circled{3}\)}
\end{problem}

\begin{answer}
C
\end{answer}

\begin{solution}
命题 \(\circled{1}\) 中，令 \(t=\frac ba\)，则 \(t>1\)，且 \(\frac ab=\frac1t<1\)，所以命题 \(\circled{1}\) 正确.

命题 \(\circled{2}\) 不正确. 例如取 \(a=1\)，\(b=2\)，\(c=3\)，\(d=6\)，则四个数均非零且 \(ad=bc=6\)，但 \(\frac ba=2\)，\(\frac cb=\frac32\)，\(\frac dc=2\)，三个相邻项之比不全相等，四个数不成等比数列，所以 \(ad=bc\) 不是四个非零实数成等比数列的充分条件.

命题 \(\circled{3}\) 的定义域为 \(x\ne0\)，且对定义域内任意 \(x\)，有 \(f(|-x|)=\log_2|-x|=\log_2|x|=f(|x|)\)，所以 \(f(|x|)\) 是偶函数. 因此正确命题的序号是 \(\circled{1}\)\(\circled{3}\)，故选 C.
\end{solution}

\begin{problem}
某生物生长过程中，在三个连续时段内的增长量都相等，在各时段内平均增长速度分别为 \(v_1\)，\(v_2\)，\(v_3\)，该生物在所讨论的整个时段内的平均增长速度为（\quad）

\choices
    {\(\frac{v_1 + v_2 + v_3}{3}\)}
    {\(\frac{\frac{1}{v_1} + \frac{1}{v_2} + \frac{1}{v_3}}{3}\)}
    {\(\sqrt[3]{v_1 v_2 v_3}\)}
    {\(\frac{3}{\frac{1}{v_1} + \frac{1}{v_2} + \frac{1}{v_3}}\)}
\end{problem}

\begin{answer}
D
\end{answer}

\begin{solution}
设三个连续时段内的增长量都为 \(s\)，则各时段的时间分别为 \(\frac{s}{v_1}\)，\(\frac{s}{v_2}\)，\(\frac{s}{v_3}\). 整个时段的总增长量为 \(3s\)，总时间为 \(\frac{s}{v_1}+\frac{s}{v_2}+\frac{s}{v_3}\)，所以整个时段内的平均增长速度为
\(\frac{3s}{\frac{s}{v_1}+\frac{s}{v_2}+\frac{s}{v_3}}=\frac{3}{\frac1{v_1}+\frac1{v_2}+\frac1{v_3}}\)，故选 D.
\end{solution}

\section{填空题}

\begin{problem}
\((1 + 2x)^5\) 的展开式中 \(x^2\) 项的系数是\fillinblank{}.
\end{problem}

\begin{answer}
\(40\)
\end{answer}

\begin{solution}
二项式定理中，\((1+2x)^5\) 的一般项为 \(\mathrm C_5^k(2x)^k\). 令 \(k=2\) 得 \(x^2\) 项，其系数为 \(\mathrm C_5^2\cdot2^2=10\cdot4=40\).
\end{solution}

\begin{problem}
已知实数 \(x\)，\(y\) 满足条件 \(\begin{cases}
x - 2y + 4 \geqslant 0,\\
3x - y - 3 \leqslant 0,\\
x \geqslant 0,\ y \geqslant 0,
\end{cases}\) 则 \(z = x + 2y\) 的最大值为\fillinblank{}.
\end{problem}

\begin{answer}
\(8\)
\end{answer}

\begin{solution}
由约束条件得 \(x\geqslant0\)，\(y\geqslant0\)，\(y\leqslant\frac{x}{2}+2\)，\(y\geqslant3x-3\). 可行域的顶点分别为 \(\left(0,0\right)\)，直线 \(x=0\) 与 \(y=\frac{x}{2}+2\) 的交点 \(\left(0,2\right)\)，直线 \(y=0\) 与 \(y=3x-3\) 的交点 \(\left(1,0\right)\)，以及由 \(\frac{x}{2}+2=3x-3\) 解得的交点 \(\left(2,3\right)\). 由于目标函数是线性函数，只需比较各顶点处的函数值：\(z\left(0,0\right)=0\)，\(z\left(0,2\right)=4\)，\(z\left(1,0\right)=1\)，\(z\left(2,3\right)=8\)，所以 \(z\) 的最大值为 \(8\).
\end{solution}

\begin{problem}
安排 \(3\) 名支教教师去 \(4\) 所学校任教，每校至多 \(2\) 人，则不同的分配方案共有\fillinblank{} 种.
\end{problem}

\begin{answer}
\(60\)
\end{answer}

\begin{solution}
每名教师均不同，且每所学校至多安排 \(2\) 人，因此分配方式只有两类. 第一类是 \(3\) 名教师分别去 \(3\) 所不同的学校，共有 \(\mathrm C_4^3\mathrm A_3^3=24\) 种；第二类是恰有 \(2\) 所学校接受教师，其中一所学校安排 \(2\) 人，另一所学校安排 \(1\) 人，共有 \(\mathrm C_4^2\mathrm C_2^1\mathrm C_3^2=36\) 种. 因此不同的分配方案共有 \(24+36=60\) 种.
\end{solution}

\begin{problem}
如图，平面内有三个向量 \(\overrightarrow{OA}\)，\(\overrightarrow{OB}\)，\(\overrightarrow{OC}\)，其中 \(\overrightarrow{OA}\) 与 \(\overrightarrow{OB}\) 的夹角为 \(120^{\circ}\)，\(\overrightarrow{OA}\) 与 \(\overrightarrow{OC}\) 的夹角为 \(30^{\circ}\)，且 \(\left|\overrightarrow{OA}\right| = \left|\overrightarrow{OB}\right| = 1\)，\(\left|\overrightarrow{OC}\right| = 2\sqrt{3}\). 若 \(\overrightarrow{OC} = \lambda \overrightarrow{OA} + \mu \overrightarrow{OB}\)（\(\lambda\)，\(\mu \in \R\)），则 \(\lambda + \mu\) 的值为\fillinblank{}.

\bitmapfigure[max width=0.34\linewidth]{img/2007/shaanxi_liberal/q16_fig1.png}
\end{problem}

\begin{answer}
\(6\)
\end{answer}

\begin{solution}
取 \(\overrightarrow{OA}\) 所在方向为 \(x\) 轴正方向，以 \(O\) 为原点. 由题图，
\(\overrightarrow{OA}=(1,0)\)，\(\overrightarrow{OB}=\left(\cos120^{\circ},\sin120^{\circ}\right)=\left(-\frac12,\frac{\sqrt3}{2}\right)\)，
\(\overrightarrow{OC}=2\sqrt3\left(\cos30^{\circ},\sin30^{\circ}\right)=\left(3,\sqrt3\right)\).

由 \(\overrightarrow{OC}=\lambda\overrightarrow{OA}+\mu\overrightarrow{OB}\)，得
\(\lambda\left(1,0\right)+\mu\left(-\frac12,\frac{\sqrt3}{2}\right)=\left(3,\sqrt3\right)\). 比较两个坐标，得到 \(\frac{\sqrt3}{2}\mu=\sqrt3\)，即 \(\mu=2\)，以及 \(\lambda-\frac12\mu=3\)，即 \(\lambda=4\). 所以 \(\lambda+\mu=6\).
\end{solution}

\section{解答题}

\begin{problem}
设函数 \(f(x) = \bs{a} \cdot \bs{b}\)，其中向量 \(\bs{a} = (m, \cos x)\)，\(\bs{b} = (1 + \sin x, 1)\)，\(x \in \R\)，且 \(f\left(\frac{\pi}{2}\right) = 2\).

\begin{enumerate}
\item 求实数 \(m\) 的值；
\item 求函数 \(f(x)\) 的最小值.
\end{enumerate}
\end{problem}

\begin{solution}
\begin{enumerate}
\item 由向量数量积的坐标运算，得 \(f(x)=m\left(1+\sin x\right)+\cos x\). 所以 \(f\left(\frac{\pi}{2}\right)=2m=2\)，从而 \(m=1\).
\item 由（1）得 \(f(x)=1+\sin x+\cos x=1+\sqrt{2}\sin\left(x+\frac{\pi}{4}\right)\). 因为 \(-1\leqslant\sin\left(x+\frac{\pi}{4}\right)\leqslant1\)，所以 \(f(x)\geqslant1-\sqrt{2}\)，且当 \(\sin\left(x+\frac{\pi}{4}\right)=-1\) 时等号成立. 因此函数 \(f(x)\) 的最小值为 \(1-\sqrt{2}\).
\end{enumerate}
\end{solution}

\begin{problem}
某项选拔共有四轮考核，每轮设有一个问题，能正确回答问题者进入下一轮考核，否则即被淘汰. 已知某选手能正确回答第一，二，三，四轮的问题的概率分别为 \(\frac{4}{5}\)，\(\frac{3}{5}\)，\(\frac{2}{5}\)，\(\frac{1}{5}\)，且各轮问题能否正确回答互不影响.

\begin{enumerate}
\item 求该选手进入第四轮才被淘汰的概率；
\item 求该选手至多进入第三轮考核的概率.
\end{enumerate}
\end{problem}

\begin{solution}
设事件 \(E_i\) 表示该选手能正确回答第 \(i\) 轮的问题，则 \(P\left(E_1\right)=\frac45\)，\(P\left(E_2\right)=\frac35\)，\(P\left(E_3\right)=\frac25\)，\(P\left(E_4\right)=\frac15\). 由各轮问题能否正确回答互不影响，以下事件的概率可以相乘.

\begin{enumerate}
\item 进入第四轮才被淘汰，表示前三轮均回答正确而第四轮回答错误，所以所求概率为 \(P\left(E_1\cap E_2\cap E_3\cap\overline{E_4}\right)=\frac45\cdot\frac35\cdot\frac25\cdot\left(1-\frac15\right)=\frac{96}{625}\).
\item 进入第四轮的必要条件是前三轮均回答正确，因此至多进入第三轮考核的事件是 \(E_1\cap E_2\cap E_3\) 的对立事件，所求概率为 \(1-P\left(E_1\cap E_2\cap E_3\right)=1-\frac45\cdot\frac35\cdot\frac25=\frac{101}{125}\).
\end{enumerate}
\end{solution}

\begin{problem}
如图，在底面为直角梯形的四棱锥 \(P - ABCD\) 中，\(AD \parallel BC\)，\(\angle ABC = 90^{\circ}\)，\(PA \perp\) 平面 \(ABCD\)，\(PA = 3\)，\(AD = 2\)，\(AB = 2\sqrt{3}\)，\(BC = 6\).

\begin{enumerate}
\item 求证：\(BD \perp\) 平面 \(PAC\)；
\item 求二面角 \(P - BD - A\) 的大小.
\end{enumerate}

\bitmapfigure[max width=0.32\linewidth]{img/2007/shaanxi_liberal/q19_fig1.png}
\end{problem}

\begin{solution}
\begin{enumerate}
\item 建立空间直角坐标系，令 \(B\) 为原点，\(BC\) 沿 \(x\) 轴正方向，\(BA\) 沿 \(y\) 轴正方向，\(PA\) 沿 \(z\) 轴正方向，则 \(A=\left(0,2\sqrt{3},0\right)\)，\(B=\left(0,0,0\right)\)，\(C=\left(6,0,0\right)\)，\(D=\left(2,2\sqrt{3},0\right)\)，\(P=\left(0,2\sqrt{3},3\right)\).
从而 \(\overrightarrow{BD}=\left(2,2\sqrt{3},0\right)\)，\(\overrightarrow{AP}=\left(0,0,3\right)\)，\(\overrightarrow{AC}=\left(6,-2\sqrt{3},0\right)\). 于是 \(\overrightarrow{BD}\cdot\overrightarrow{AP}=0\)，且 \(\overrightarrow{BD}\cdot\overrightarrow{AC}=12-12=0\)，所以 \(\overrightarrow{BD}\) 是平面 \(PAC\) 的法向方向；为确认直线与平面相交，由 \(\overrightarrow{AP}\times\overrightarrow{AC}\) 可得平面 \(PAC\) 的方程为 \(x+\sqrt{3}y-6=0\)；直线 \(BD\) 上的点为 \(\left(2t,2\sqrt{3}t,0\right)\)，取 \(t=\frac34\) 即满足该平面方程；因此 \(BD\perp\) 平面 \(PAC\).
\item 平面 \(ABD\) 的一个法向量为 \(\bs{n}_1=\left(0,0,1\right)\)，平面 \(PBD\) 的一个法向量为
\[
\bs{n}_2=\overrightarrow{BD}\times\overrightarrow{BP}=\left(2,2\sqrt{3},0\right)\times\left(0,2\sqrt{3},3\right)=\left(6\sqrt{3},-6,4\sqrt{3}\right)
\]
因此两平面的夹角 \(\theta\) 满足 \(\cos\theta=\frac{|\bs{n}_1\cdot\bs{n}_2|}{|\bs{n}_1||\bs{n}_2|}=\frac{4\sqrt{3}}{\sqrt{\left(6\sqrt{3}\right)^2+\left(-6\right)^2+\left(4\sqrt{3}\right)^2}}=\frac{4\sqrt{3}}{8\sqrt{3}}=\frac12\)，故二面角 \(P-BD-A\) 的大小为 \(\theta=60^{\circ}\).
\end{enumerate}
\end{solution}

\begin{problem}
已知实数列 \(\{a_n\}\) 是等比数列，其中 \(a_7 = 1\)，且 \(a_4\)，\(a_5 + 1\)，\(a_6\) 成等差数列.

\begin{enumerate}
\item 求数列 \(\{a_n\}\) 的通项公式；
\item 数列 \(\{a_n\}\) 的前 \(n\) 项和记为 \(S_n\)，证明：\(S_n < 128\)（\(n = 1\)，\(2\)，\(3\)，\(\cdots\)）.
\end{enumerate}
\end{problem}

\begin{solution}
\begin{enumerate}
\item 设等比数列 \(\{a_n\}\) 的公比为 \(q\). 由于 \(a_7=1\neq0\)，所以 \(q\neq0\)，并且 \(a_4=q^{-3}\)，\(a_5=q^{-2}\)，\(a_6=q^{-1}\). 由 \(a_4\)，\(a_5+1\)，\(a_6\) 成等差数列，得 \(2\left(q^{-2}+1\right)=q^{-3}+q^{-1}\). 两边同乘 \(q^3\)，得 \(2q^3-q^2+2q-1=0\)，即 \(\left(2q-1\right)\left(q^2+1\right)=0\). 因为 \(q\) 为实数，故 \(q=\frac12\). 再由 \(a_7=a_1q^6=1\)，得 \(a_1=64\)，所以数列的通项公式为 \(a_n=64\left(\frac12\right)^{n-1}=2^{7-n}\).
\item 由等比数列前 \(n\) 项和公式，得 \(S_n=\frac{64\left(1-\left(\frac12\right)^n\right)}{1-\frac12}=128\left(1-2^{-n}\right)=128-2^{7-n}\). 当 \(n=1\)，\(2\)，\(3\)，\(\cdots\) 时，\(2^{7-n}>0\)，所以 \(S_n<128\).
\end{enumerate}
\end{solution}

\begin{problem}
已知函数 \(f(x) = a x^3 + b x^2 + c x\) 在区间 \([0, 1]\) 上是增函数，在区间 \((-\infty, 0)\)，\((1, +\infty)\) 上是减函数. 又 \(f'\left(\frac{1}{2}\right) = \frac{3}{2}\).

\begin{enumerate}
\item 求 \(f(x)\) 的解析式；
\item 若在区间 \([0, m]\)（\(m > 0\)）上恒有 \(f(x) \leqslant x\) 成立，求 \(m\) 的取值范围.
\end{enumerate}
\end{problem}

\begin{solution}
\begin{enumerate}
\item 由于 \(f\) 在 \((-\infty,0)\) 上递减且可导，对 \(x<0\) 有 \(f'(x)\leqslant0\)；由于 \(f\) 在 \((0,1)\) 上递增，对 \(0<x<1\) 有 \(f'(x)\geqslant0\). 又 \(f'\) 连续，所以 \(f'(0)\) 同时受左侧的非正性和右侧的非负性约束，从而 \(f'(0)=0\). 在 \(1\) 的左侧，\(f\) 在 \((0,1)\) 上递增，故 \(f'(x)\geqslant0\)（\(0<x<1\)）；在 \(1\) 的右侧，\(f\) 在 \((1,+\infty)\) 上递减，故 \(f'(x)\leqslant0\)（\(x>1\)）. 由 \(f'\) 在 \(1\) 处连续，\(f'(1)\) 同时受左侧的非负性和右侧的非正性约束，因此 \(f'(1)=0\). 又 \(f'(x)=3ax^2+2bx+c\)，所以 \(f'(x)=k x\left(x-1\right)\). 由 \(f'\left(\frac12\right)=\frac32\)，得 \(-\frac{k}{4}=\frac32\)，从而 \(k=-6\). 比较系数可得 \(3a=-6\)，\(2b=6\)，\(c=0\)，故 \(a=-2\)，\(b=3\)，\(c=0\)，即 \(f(x)=-2x^3+3x^2\).
\item 由（1）得 \(x-f(x)=x+2x^3-3x^2=x\left(2x-1\right)\left(x-1\right)\). 当 \(0\leqslant x\leqslant\frac12\) 时，\(x\geqslant0\)，\(2x-1\leqslant0\)，\(x-1\leqslant0\)，所以 \(x-f(x)\geqslant0\)，即 \(f(x)\leqslant x\). 当 \(\frac12<x<1\) 时，\(x-f(x)<0\)，即 \(f(x)>x\). 因此要使 \([0,m]\) 上恒有 \(f(x)\leqslant x\)，必须且只需 \(0<m\leqslant\frac12\).
\end{enumerate}
\end{solution}

\begin{problem}
已知椭圆 \(C: \frac{x^2}{a^2} + \frac{y^2}{b^2} = 1\)（\(a > b > 0\)）的离心率为 \(\frac{\sqrt{6}}{3}\)，短轴一个端点到右焦点的距离为 \(\sqrt{3}\).

\begin{enumerate}
\item 求椭圆 \(C\) 的方程；
\item 设直线 \(l\) 与椭圆 \(C\) 交于 \(A\)，\(B\) 两点，坐标原点 \(O\) 到直线 \(l\) 的距离为 \(\frac{\sqrt{3}}{2}\)，求 \(\triangle AOB\) 面积的最大值.
\end{enumerate}
\end{problem}

\begin{solution}
\begin{enumerate}
\item 设椭圆的半焦距为 \(c\)，则 \(c^2=a^2-b^2\)，且由离心率条件得 \(\frac{c}{a}=\frac{\sqrt{6}}{3}\)，所以 \(c^2=\frac23a^2\)，\(b^2=\frac13a^2\). 取短轴端点 \(\left(0,b\right)\) 和右焦点 \(\left(c,0\right)\)，由题意得 \(c^2+b^2=3\)，从而 \(a^2=3\)，\(b^2=1\). 故椭圆 \(C\) 的方程为 \(\frac{x^2}{3}+y^2=1\).
\item 设直线 \(l\) 的单位法向量为 \(\bs{n}=\left(\cos\theta,\sin\theta\right)\)，并记 \(d=\frac{\sqrt3}{2}\)，则 \(l\) 可表示为 \(x\cos\theta+y\sin\theta=d\). 取直线的单位方向向量 \(\bs{v}=\left(-\sin\theta,\cos\theta\right)\)，直线上任一点可表示为 \(\left(x,y\right)=d\bs{n}+t\bs{v}\). 代入椭圆方程，得到关于 \(t\) 的二次方程
\[
K t^2+\frac{4d}{3}\sin\theta\cos\theta\,t+d^2\left(\frac{\cos^2\theta}{3}+\sin^2\theta\right)-1=0
\]
其中 \(K=\frac{\sin^2\theta}{3}+\cos^2\theta\). 该二次方程的判别式为
\[
\Delta=\left(\frac{4d}{3}\sin\theta\cos\theta\right)^2-4K\left[d^2\left(\frac{\cos^2\theta}{3}+\sin^2\theta\right)-1\right]=4\left(K-\frac{d^2}{3}\right)=4K-1
\]
其中恒等式可直接展开验证：令 \(u=\sin^2\theta\)，\(v=\cos^2\theta\)，则 \(u+v=1\)，且
\[
K\left(\frac{\cos^2\theta}{3}+\sin^2\theta\right)-\frac49\sin^2\theta\cos^2\theta
=\left(\frac{u}{3}+v\right)\left(\frac{v}{3}+u\right)-\frac49uv
=\frac{u^2+v^2+2uv}{3}=\frac{(u+v)^2}{3}=\frac13
\]
因为 \(\bs{v}\) 是单位向量，两个根之差就是弦 \(AB\) 的长度，所以 \(|AB|=\frac{\sqrt{\Delta}}{K}=\frac{\sqrt{4K-1}}{K}\). 又 \(K=\frac13+\frac23\cos^2\theta\)，故 \(\frac13\leqslant K\leqslant1\).

三角形 \(AOB\) 的高为原点 \(O\) 到直线 \(l\) 的距离 \(d\)，所以 \(S_{\triangle AOB}=\frac12d|AB|\). 于是
\[
S_{\triangle AOB}^2=\frac{d^2}{4}\cdot\frac{4K-1}{K^2}=\frac{3\left(4K-1\right)}{16K^2}
\]
又 \(\frac{4K-1}{K^2}=\frac4K-\frac1{K^2}\)，且
\[
\frac{\mathrm d}{\mathrm dK}\left(\frac{4K-1}{K^2}\right)=\frac{2-4K}{K^3}
\]
由于 \(K>0\)，该函数在 \(K<\frac12\) 时递增，在 \(K>\frac12\) 时递减，所以当 \(K=\frac12\) 时取得最大值. 此时 \(\cos^2\theta=\frac14\)，该取值可以实现，且 \(|AB|=2\)，从而三角形 \(AOB\) 面积的最大值为 \(\frac12\cdot\frac{\sqrt3}{2}\cdot2=\frac{\sqrt3}{2}\).
\end{enumerate}
\end{solution}
